% SI derivation text (session P0-W3.4; contract v3, thresholds block v3, ignorance block v2).
% \input-able into the Supplementary Information in place of theory/derivation_body.tex, under the SI's
% own \section{Full derivation}: this file starts at \subsection level and does not number sections.
% Requirements of the including document: amsmath, array, and \input of theory/macros.tex.
% Labels carry the prefix thsi: so they collide neither with the paper's labels nor with th: of
% theory/derivation_body.tex.
% Symbolic only: no literature value appears here. The few registered numbers reach the text only
% through the generated macros \NumPcred, \NumPcredSweepLow, \NumPcredSweepHigh, \NumMexpt and
% \NumMwindow (W8 generates them; see ops/reports/P0-W3.4.md). Until they exist each renders a
% visible marker, and a release build stops on it.
\providecommand{\thsiNum}[1]{\ifdefined\nacrePending\nacrePending{macro #1}\else\texttt{[#1]}\fi}
\providecommand{\NumPcred}{\thsiNum{NumPcred}}
\providecommand{\NumPcredSweepLow}{\thsiNum{NumPcredSweepLow}}
\providecommand{\NumPcredSweepHigh}{\thsiNum{NumPcredSweepHigh}}
\providecommand{\NumMexpt}{\thsiNum{NumMexpt}}
\providecommand{\NumMwindow}{\thsiNum{NumMwindow}}
\providecommand{\Apatch}{A_\mathrm{patch}}
\providecommand{\Nreq}{N_\mathrm{req}}
\providecommand{\Nnuc}{N_\mathrm{nuc}}
\providecommand{\mexpt}{m_\mathrm{expt}}
\providecommand{\Ecmax}{E_\mathrm{cell}^{\max}}
\providecommand{\Ecmin}{E_\mathrm{cell}^{\min}}

This section derives every relation the screen evaluates. Each part opens with the physical
argument; the algebra follows it. The notation is that of the symbol register of the theory package
(listed in the parameter tables), and every displayed relation carries its SI units in brackets. Two
regimes are kept apart throughout: the water-processed, alkaline electrode slurry at room temperature
(\Rs) and the non-aqueous, fluorine-free cell with the Al collector at the cathode's potential (\Rc).
A relation is used only in the regime whose reference states it assumes. The links C1--C5 name what
the proposed mechanism must achieve: anchor the template (C1), bind Ca$^{2+}$ (C2), nucleate the
mineral (C3), form a continuous film (C4) and keep it stable (C5). Meeting the conditions below
makes the mechanism thermodynamically possible under the stated framing; it does not show that a
barrier works, and no rate, capacity or device statement follows from them.

%----------------------------------------------------------------------------------------------
\subsection{Adhesion of the template in water (C1)}\label{thsi:sec:adhesion}

\textit{Physics.} In the slurry the template never meets a dry oxide. The hydroxylated oxide and
the polymer are both covered by water, and a contact between them forms only by displacing that
water. The competing states are therefore the contact and the two water-wetted surfaces, not the
contact and two bare surfaces.

The Dupr\'e work to separate a contact between the oxide (o) and the polymer (p) in water (w)
creates an oxide--water and a polymer--water interface and destroys the oxide--polymer one:
\begin{equation}\label{thsi:eq:dupre}
\begin{gathered}
\Wwet=\gow+\gpw-\gop=\Wdry+2\gw-W_\mathrm{ow}-W_\mathrm{pw},\\
\Wdry=\go+\gp-\gop \un{J\,m^{-2}},
\end{gathered}
\end{equation}
where $W_\mathrm{ow}=\go+\gw-\gow$ and $W_\mathrm{pw}=\gp+\gw-\gpw$ are the works with which water
binds each surface. A large dry adhesion is offset by those works, so $\Wwet$ can be small or
negative for a hydrophilic oxide and polymer. The bare oxide has $\Wwet=0$ and fixes the zero.
The route that decides C1 is the potential of mean force of a template fragment in explicit water,
$\Wwet=\Delta F_\mathrm{PMF}/A_\mathrm{c}$, which contains the displacement penalty and the entropy
of the released water; it needs a PMF that plateaus at large separation and one contact-area
definition for all fragments.

\paragraph{C1 inputs in P0.} No PMF exists for P0, so the registered C1 verdict is a gap and counts
as inconclusive. Two estimates are reported beside it. (i) A flagged analog: works of adhesion in
water measured for catechol-bearing films on other oxides. A pull-off measurement records a work
only for contacts that adhere, so the analog has positive support and cannot represent displacement
by water; its verdict is conditional on the analog. (ii) The Dupr\'e form with ignorance-bounded
energies, $\gow\in[0,\go+\gw]$ and $\gop\in[0,\go+\gp]$ (Section~\ref{thsi:sec:branches}), and
$\gpw$ the sourced template--water energy truncated at zero. It places $\Wwet$ in
$[\gpw-\go-\gp,\ \gpw+\go+\gw]$. Since $\gpw\le\gp+\gw$, the lower end is at most $\gw-\go$, negative
whenever the dry oxide's surface energy exceeds water's: thermodynamic bounds alone admit both signs
of the wet adhesion, which is why a PMF is needed. Masses under these bounds are properties of the
prior, never a finding; the flip value $\gop^{*}=\gow+\gpw-m_W$ is reported instead.

%----------------------------------------------------------------------------------------------
\subsection{The driving force: uniform chemical potential}\label{thsi:sec:uniform}

\textit{Physics.} A nucleus is built from ions taken out of the solution, and it gains, per formula
unit, the chemical potential the ions carry minus that of the solid. The template sits in an open
interfacial region that exchanges ions with the slurry. At equilibrium nothing is gained by moving
an ion between bulk and layer, so every ion carries the same electrochemical potential everywhere.
The template can concentrate ions; it cannot make them carry more free energy.

\paragraph{One definition.} With $\Ksp$ the solubility product of the registered framing (calcite
and amorphous CaCO$_3$ are both registered, each run as a whole chain with its own $\Ksp$, $\Vm$ and
$\gnl$),
\begin{equation}\label{thsi:eq:omega}
\Om=\frac{\aCa\,\aCO}{\Ksp} \un{1},\qquad \Dmu=\kB T\ln\Om \un{J},
\end{equation}
per CaCO$_3$ formula unit, with molal activities. Nucleation requires $\Om>1$. For a hydrated
framing the ion activity product carries $a_\mathrm{w}^{n}$; keeping $\aCa\aCO$ overstates $\Om$ by
$a_\mathrm{w}^{-n}$ at a drying state ($a_\mathrm{w}<1$), so drying states are optimistic, $a_\mathrm{w}$
is reported there, and the hydrated ratio is a labelled sensitivity.

\paragraph{Equilibrium of the open region.} The electrochemical potential of ion $i$ at $\mathbf r$ is
\begin{equation}\label{thsi:eq:echem}
\tilde\mu_i(\mathbf r)=\mu_i^{\circ}+w_i^\mathrm{sp}(\mathbf r)+\kB T\ln\frac{\gamma_i(\mathbf r)m_i(\mathbf r)}{m^{\circ}}+z_ie\psi(\mathbf r) \un{J},
\end{equation}
with $w_i^\mathrm{sp}$ the specific interaction with the template, $\gamma_i$ the non-specific
activity coefficient and $\psi$ the mean potential (both $w^\mathrm{sp}$ and $\psi$ vanish in bulk).
Equilibrium with the slurry reservoir is $\tilde\mu_i(\mathbf r)=\tilde\mu_{i,\mathrm{bulk}}$ for every
ion and every $\mathbf r$. Solving for $m_i$ gives the local enrichment
$m_i/m_{i,\mathrm{bulk}}=(\gamma_{i,\mathrm{bulk}}/\gamma_i)\exp[-(w_i^\mathrm{sp}+z_ie\psi)/\kB T]$: the
Boltzmann enrichment is a consequence of uniform potentials, not an exception to them. The driving
force for adding a formula unit at $\mathbf r$ is then
\begin{equation}\label{thsi:eq:dmu-uniform}
\Dmu(\mathbf r)=\tilde\mu_\mathrm{Ca}(\mathbf r)+\tilde\mu_{\mathrm{CO_3}}(\mathbf r)-\mu_\mathrm{CaCO_3(s)}=\kB T\ln\Omb \un{J},
\end{equation}
the same at a binding site, in the diffuse layer and in bulk.

\paragraph{Specific binding.} At a site, binding lowers the standard part of $\tilde\mu_\mathrm{Ca}$ by
$-\DGseq/\NA$ and raises the concentration term by $\kB T\ln\Eenr=-\DGseq/\NA$. The two cancel:
\begin{equation}\label{thsi:eq:no-gain}
\Dmu_\mathrm{site}-\Dmu_\mathrm{bulk}=\kB T\ln\Eenr+\DGseq/\NA=0 \un{J}.
\end{equation}
\paragraph{A mean potential.} At an anionic template Ca$^{2+}$ is enriched by $\e^{-2e\psi/\kB T}$ and
CO$_3^{2-}$ depleted by $\e^{+2e\psi/\kB T}$. The ion product is unchanged, and in
Eq.~\eqref{thsi:eq:dmu-uniform} the potential enters as $(z_\mathrm{Ca}+z_{\mathrm{CO_3}})e\psi=0$
because the solid is neutral.

\paragraph{Result.} The nucleus at the template is driven by $\kB T\ln\Omb$ at the regime activities.
The template changes the barrier only through interfacial energies. A local driving force above the
bulk one needs a non-equilibrium state with a time scale (drying faster than exchange, a layer
loaded at another activity, a local pH excursion); P0 makes no statement about such states.

%----------------------------------------------------------------------------------------------
\subsection{Carbonate speciation and the activity model}\label{thsi:sec:speciation}

\textit{Physics.} $\Om$ needs the activities of the free ions. In an alkaline, carbonate-bearing
slurry part of the calcium and carbon is bound in ion pairs and part of the carbon is bicarbonate or
dissolved CO$_2$, so the totals overstate the free ions. The registered state fixes pH, the dissolved
inorganic carbon $C_\mathrm{T}$, the total dissolved calcium $\mathrm{Ca}_\mathrm{T}$ and the ionic
strength $I$, in a closed solution that has not yet precipitated.

With $a_\mathrm{H}=10^{-\mathrm{pH}}$ and $y\equiv\aCO$, the thermodynamic constants (at $I=0$) tie
every species to $y$ and $\aCa$: $a_{\mathrm{HCO_3}}=10^{\mathrm{p}K_2-\mathrm{pH}}y$,
$a_{\mathrm{CO_2}}=10^{\mathrm{p}K_1+\mathrm{p}K_2-2\mathrm{pH}}y$,
$a_{\mathrm{CaCO_3^0}}=K_{\mathrm{CaCO_3^0}}\aCa y$,
$a_{\mathrm{CaHCO_3^+}}=K_{\mathrm{CaHCO_3^+}}\aCa a_{\mathrm{HCO_3}}$ and
$a_{\mathrm{CaOH^+}}={}^{*}K_{\mathrm{CaOH^+}}\aCa a_\mathrm{w}/a_\mathrm{H}$. Dividing each activity by
its coefficient ($\gamma_{\pm1}$ or $\gamma_{\pm2}$ by charge; one for neutral species) and summing gives
the carbon and calcium balances in units of the standard molality:
\begin{equation}\label{thsi:eq:balance}
m_{C_\mathrm{T}}=A_\mathrm{c}y+P_\mathrm{c}\aCa y,\qquad m_{\mathrm{Ca}_\mathrm{T}}=D_\mathrm{c}\aCa+P_\mathrm{c}\aCa y \un{1},
\end{equation}
\begin{equation}\label{thsi:eq:factors}
\begin{gathered}
A_\mathrm{c}=10^{\mathrm{p}K_1+\mathrm{p}K_2-2\mathrm{pH}}+\frac{10^{\mathrm{p}K_2-\mathrm{pH}}}{\gamma_{\pm1}}+\frac{1}{\gamma_{\pm2}},\\
P_\mathrm{c}=K_{\mathrm{CaCO_3^0}}+\frac{K_{\mathrm{CaHCO_3^+}}10^{\mathrm{p}K_2-\mathrm{pH}}}{\gamma_{\pm1}},\\
D_\mathrm{c}=\frac{1}{\gamma_{\pm2}}+\frac{{}^{*}K_{\mathrm{CaOH^+}}a_\mathrm{w}10^{\mathrm{pH}}}{\gamma_{\pm1}} \un{1},
\end{gathered}
\end{equation}
with $m_{C_\mathrm{T}}=C_\mathrm{T}/(\rho_\mathrm{w}m^{\circ})$ and
$m_{\mathrm{Ca}_\mathrm{T}}=\mathrm{Ca}_\mathrm{T}/(\rho_\mathrm{w}m^{\circ})$ in dilute solution.
Eliminating $\aCa=m_{\mathrm{Ca}_\mathrm{T}}/(D_\mathrm{c}+P_\mathrm{c}y)$ leaves
$A_\mathrm{c}P_\mathrm{c}y^{2}+By-m_{C_\mathrm{T}}D_\mathrm{c}=0$ with
$B=A_\mathrm{c}D_\mathrm{c}+P_\mathrm{c}(m_{\mathrm{Ca}_\mathrm{T}}-m_{C_\mathrm{T}})$. The product of the
roots is negative, so exactly one root is positive; in the form that stays exact as $P_\mathrm{c}\to0$,
\begin{equation}\label{thsi:eq:aco3}
\aCO=\frac{2m_{C_\mathrm{T}}D_\mathrm{c}}{B+\sqrt{B^{2}+4A_\mathrm{c}P_\mathrm{c}m_{C_\mathrm{T}}D_\mathrm{c}}},\qquad
\aCa=\frac{m_{\mathrm{Ca}_\mathrm{T}}}{D_\mathrm{c}+P_\mathrm{c}\aCO} \un{1}.
\end{equation}
With no ion pairs or no calcium, $\aCO=m_{C_\mathrm{T}}/A_\mathrm{c}$, the carbonate system alone; pairs
can only lower it. The free Ca$^{2+}$ concentration $\cb=\aCa\rho_\mathrm{w}m^{\circ}/\gamma_{\pm2}$ is the
$\cb$ of the binding relations below: the free ion, not the total.

\paragraph{Activity model.} Davies:
$\log_{10}\gamma_{\pm z}=-A_\mathrm{DH}z^{2}\big[s/(1+s)-b_\mathrm{D}s^{2}\big]$ with $s=\sqrt{I/m^{\circ}}$, so
$\gamma_{\pm2}=\gamma_{\pm1}^{4}$ and every coefficient tends to one as $I\to0$. The slope
$\dd\log_{10}\gamma/\dd s$ changes sign at the positive root $s_*$ of $2b_\mathrm{D}s(1+s)^{2}=1$; beyond
$I_*=m^{\circ}s_*^{2}$ the ionic-strength dependence is carried by the empirical linear term alone.
The registered validity bound therefore satisfies $I_\mathrm{val}\le I_*$, and a draw with
$I>I_\mathrm{val}$ is refused and counted, never extrapolated. Neutral species take coefficient one,
and $I$ is the registered input rather than recomputed from the speciation; both are stated
limitations.

%----------------------------------------------------------------------------------------------
\subsection{Ca$^{2+}$ binding at the template (C2, a descriptor)}\label{thsi:sec:binding}

\textit{Physics.} The template can raise the local Ca$^{2+}$ concentration only if it holds Ca$^{2+}$
more strongly than water does. The free energy must start from hydrated Ca$^{2+}$ at the regime
activity, pay for preparing the site at the regime pH, and end in a well-defined local state.

For an acid site L--H that binds only when deprotonated, the cycle (deprotonate at the regime pH,
then transfer Ca$^{2+}$ from bulk to the site) gives
\begin{equation}\label{thsi:eq:cycle}
\begin{gathered}
\DGseq=\DGsite+\Rg T\ln\big(1+10^{\,\mathrm{p}K_\mathrm{a}-\mathrm{pH}}\big),\\
\DGsite\equiv-\Rg T\ln\Big[\frac{1}{\Vsite}\int_\mathrm{site}\e^{-w(\mathbf r)/\kB T}\dd^{3}r\Big] \un{J\,mol^{-1}},
\end{gathered}
\end{equation}
with $w$ the PMF of one ion relative to bulk and $\Vsite$ the binding-layer volume per site, fixed
once. A binding constant to a standard concentration $c^{\circ}$ compares with $\DGsite$ only after
$\Delta G^{\circ}_\mathrm{bind}=\DGsite-\Rg T\ln(\NA c^{\circ}\Vsite)$. In the dilute limit with
non-interacting sites the enrichment is the concentration ratio
$\Eenr=\cl/\cb=\exp(-\DGseq/\Rg T)$; by Section~\ref{thsi:sec:uniform} it never enters $\Om$. For
$\Gs$ sites per area in a layer $\ds$ thick (Langmuir),
\begin{equation}\label{thsi:eq:langmuir}
\begin{gathered}
x=\frac{\Eenr\cb}{\cmax},\quad \cmax=\frac{\Gs}{\ds},\quad \Theta=\frac{x}{1+x},\\
\cl=\cmax\Theta\le\min(\Eenr\cb,\cmax),\\
\GCa=\Gs\Theta-\ds\cb \un{1;\ mol\,m^{-2}},
\end{gathered}
\end{equation}
where $\GCa$ is the Gibbs excess: the bound amount per area less what the layer volume holds at the
bulk concentration. A strongly binding but sparse template is limited by its site density. These
quantities, and the adsorption term of the next part, are what C2 reports. C2 does not gate: its only
route to nucleation is the template--liquid energy, which C3 already judges. The mechanism's premise,
$\DGseq<0$, is reported as a finding (it holds, fails or is undecided under the shared decision rule).

\paragraph{Geometric bounds on $\Gs$ and $\ds$.} Each site holds one ion, and two bound ions cannot
be closer than their contact distance $2r_\mathrm{Ca}$. The densest planar packing of equal discs
gives each the area $2\sqrt3\,r^{2}$, so the footprint of a site obeys
$a_\mathrm{site}\ge2\sqrt3\,r_\mathrm{Ca}^{2}$ (one interfacial layer of sites). A volume of template of
mass density $\rho_\mathrm{PDA}$ holds at most $\rho_\mathrm{PDA}/M_\mathrm{u}$ moles of coordinating
groups per volume, with $M_\mathrm{u}$ the least template mass per mole of groups. Hence
\begin{equation}\label{thsi:eq:gamma-bound}
\begin{gathered}
0<\Gs\le\min\Big(\frac{1}{2\sqrt3\,\NA r_\mathrm{Ca}^{2}},\ \frac{\rho_\mathrm{PDA}d_\mathrm{layer}}{M_\mathrm{u}}\Big),\\
\max\big(2r_\mathrm{Ca},\ \Gs\Vm\big)\le\ds\le d_\mathrm{layer} \un{mol\,m^{-2};\ m}.
\end{gathered}
\end{equation}
The layer must contain the ion ($2r_\mathrm{Ca}$), cannot hold Ca$^{2+}$ more densely than a solid Ca
phase ($\Gs/\ds\le1/\Vm$) and lies inside the template layer. $r_\mathrm{Ca}$ is the ionic radius at
the lowest registered coordination number, the smallest admissible, so the bound holds at any higher
coordination. The Ca--O distance of the hydrated ion exceeds $r_\mathrm{Ca}$ and would tighten both
bounds beyond what geometry allows, so it is used only as a labelled sensitivity; while
$r_\mathrm{Ca}$ is unsourced the $2r_\mathrm{Ca}$ term is dropped, which keeps the bound valid.

%----------------------------------------------------------------------------------------------
\subsection{Where the bound Ca$^{2+}$ goes: the Gibbs--Langmuir term}\label{thsi:sec:gibbs}

\textit{Physics.} Since the driving force is the bulk one, the only places left for the template's
binding are the interfacial energies. Bound ions are an interfacial excess, and the Gibbs adsorption
equation says how such an excess changes the energy of its interface.

\paragraph{Template--liquid interface.} At fixed $T$, $\dd\gsl=-\sum_i\Gamma_i\dd\mu_i$. Taking the
adsorbing component as Ca$^{2+}$ with its compensating charge (a neutral solute) and raising the
Ca$^{2+}$ activity from zero (the Ca-free template, $\gslz$) to the regime value, with the excess of
Eq.~\eqref{thsi:eq:langmuir} and $\dd\mu_\mathrm{Ca}=\Rg T\dd\ln c$,
\begin{equation}\label{thsi:eq:gibbs}
\begin{aligned}
\dga&\equiv\gsl-\gslz=-\int_0^{\cb}\GCa(c)\,\Rg T\frac{\dd c}{c}\\
&=-\Rg T\big[\Gs\ln(1+x)-\ds\cb\big] \un{J\,m^{-2}},
\end{aligned}
\end{equation}
because $x(c)$ is linear in $c$. Its properties: $\partial\dga/\partial\DGseq=\Gs\Theta>0$, so deeper
binding lowers $\gsl$; $\partial\dga/\partial\mu_\mathrm{Ca}=-\GCa$; a repelling template has
$\dga\to\Rg T\ds\cb>0$; at strong binding
$\dga\to\Gs\DGseq+\Rg T\Gs\ln(\cmax/\cb)+\Rg T\ds\cb$, whose last term survives but is small when
$\ds\cb\ll\Gs$. If $\gsl$ is computed at the regime activities, $\dga$ is already inside it and is not
added.

\paragraph{The neutral-solute approximation and the double-layer limitation.} The cancellation of
the mean potential in Section~\ref{thsi:sec:uniform} concerns the ion product, not $\gsl$. Binding
Ca$^{2+}$ to anionic sites changes the surface charge density $\sigma$ and with it the diffuse-layer
free energy, which is part of $\gsl$; for an electrified interface the Gibbs equation carries a
Lippmann term $-\sigma\dd\psi_0$. Equation~\eqref{thsi:eq:gibbs} omits it. In the low-potential
Gouy--Chapman limit the diffuse layer holds $\sigma^{2}/(2\varepsilon_\mathrm{r}\varepsilon_0\kappa_\mathrm{D})$
per area, with $\kappa_\mathrm{D}^{2}\propto I$; at a fixed change of charge the omitted term falls as
$I^{-1/2}$ and vanishes when binding leaves the charge unchanged. It is an unquantified correction of
unassigned sign, largest at low ionic strength for a template whose charge binding changes strongly;
an in-situ $\gsl$ with explicit ions contains it.

\paragraph{Template--nucleus interface.} $\gsn$ is a grand-canonical excess at the regime chemical
potentials as well. If the contact exchanges Ca$^{2+}$ with the reservoir, any Ca$^{2+}$ excess
$\GCasn$ at the contact (an ion bridge between a carboxylate and the crystal, say) is inside it, and
\begin{equation}\label{thsi:eq:gibbs-sn}
\frac{\partial\gsn}{\partial\mu_\mathrm{Ca}}=-\GCasn,\qquad
\frac{\partial(\gsl-\gsn)}{\partial\mu_\mathrm{Ca}}=-(\GCasl-\GCasn) \un{mol\,m^{-2}},
\end{equation}
with $\GCasl\equiv\GCa$. A contact sealed from the solution holds a fixed amount instead, and the
relation does not apply. The screen takes $\gsn$ from contact data at the regime state and never adds
$\DGseq$ or $\dga$ to it; this is a rule for the inputs, not a claim that $\GCasn=0$.

%----------------------------------------------------------------------------------------------
\subsection{Nucleation on the template and the sign convention}\label{thsi:sec:nucleation}

\textit{Physics.} A cluster gains $\Dmu$ per formula unit but pays interfacial free energy for every
unit of area it creates; the barrier is the top of that trade-off. On a substrate the cluster
replaces part of its surface with a cheaper substrate contact, which lowers the barrier without
changing the driving force.

A sphere of radius $r$ with molecular volume $\vm=\Vm/\NA$ has
$\Delta G_\mathrm{hom}(r)=-\tfrac{4\pi r^{3}}{3}\Dmu/\vm+4\pi r^{2}\gnl$, so
\begin{equation}\label{thsi:eq:hom}
r^{*}=\frac{2\gnl\vm}{\Dmu} \un{m},\qquad \DGhom=\frac{16\pi\gnl^{3}\vm^{2}}{3\Dmu^{2}} \un{J}.
\end{equation}
A spherical cap of contact angle $\theta$ (measured through the nucleus) on a flat substrate has
$\Delta G_\mathrm{het}(r)=-V_\mathrm{cap}\Dmu/\vm+\gnl A_\mathrm{cap}+(\gsn-\gsl)A_\mathrm{base}$.
Stationarity in $\theta$ at fixed volume is Young's relation, and substituting it factorises the
whole free energy:
\begin{equation}\label{thsi:eq:young}
\begin{gathered}
\cos\theta=\frac{\gsl-\gsn}{\gnl},\qquad \gsn=\gsnz+\dgw;\\
\DGhet=f(\theta)\DGhom,\qquad f(\theta)=\frac{(2+\cos\theta)(1-\cos\theta)^{2}}{4} \un{1;\ J}.
\end{gathered}
\end{equation}
$f(0)=0$, $f(\pi/2)=\tfrac12$, $f(\pi)=1$ and $f'=\tfrac34\sin^{3}\theta\ge0$, so
$0\le\DGhet\le\DGhom$. The factorisation needs one and the same $\gnl$ on the cap surface and in
Young, so the bare $\gnl$ stays in the denominator.

\paragraph{The water-displacement term.} Displacing the template's bound water under the contact disc
costs free energy over the contact only: $\dgw\equiv\gsn-\gsnz>0$ is a cost, placed on the substrate
side. Then $\partial f/\partial\dgw=3\sin^{2}\theta/(4\gnl)\ge0$: a cost raises the barrier. Writing
$\geff=\gnl-\dgw$ in the denominator instead would give
$\partial\cos\theta/\partial\dgw=(\gsl-\gsnz)/\geff^{2}>0$ on the wetting side, so a cost would lower
the barrier, and it would break the factorisation. If $\gsn$ is a full in-water energy, $\dgw\equiv0$.

\paragraph{Consequence of binding for the barrier.} With $\gsl=\gslz+\dga$,
\begin{equation}\label{thsi:eq:young-ads}
\begin{gathered}
\cos\theta=\frac{\gslz+\dga-\gsnz-\dgw}{\gnl},\\
\frac{\partial\cos\theta}{\partial\DGseq}=\frac{\Gs\Theta}{\gnl}>0 \un{1;\ mol\,J^{-1}}.
\end{gathered}
\end{equation}
At fixed $\gslz$ and $\gsn$, stronger binding lowers $\cos\theta$ and raises the barrier: ions held at
the template--liquid interface must be stripped from the area the nucleus covers, the physics of an
adsorbed inhibitor. Binding promotes nucleation only if $\gsn$ falls by more than $\gsl$, that is, by
Eq.~\eqref{thsi:eq:gibbs-sn}, only if the contact holds more Ca$^{2+}$ excess than the template--liquid
interface. Lowering $\gnl$ lowers $f$ only when $\gsl>\gsn$.

%----------------------------------------------------------------------------------------------
\subsection{No double counting}\label{thsi:sec:nodouble}

\textit{Physics.} The template is both an interface and a ligand field, so its binding could be
counted twice: once as a lower interfacial energy and once as a higher driving force. The open system
at fixed bulk chemical potentials has one free energy, and writing it down shows where each effect
sits.

Forming a cap of radius $r$ transfers $n=V_\mathrm{cap}/\vm$ formula units from the reservoir at the
bulk chemical potentials, creates the cap surface and replaces template--liquid by template--nucleus
interface over the base. The grand potential change is
\begin{equation}\label{thsi:eq:grand}
\Delta\Omega_\mathrm{G}(r)=-n\,\kB T\ln\Omb+\gnl A_\mathrm{cap}+(\gsn-\gsl)A_\mathrm{base} \un{J},
\end{equation}
with grand-canonical interfacial energies at the bulk chemical potentials. Within capillarity this
is exact; the reservoir term carries no template dependence, and by Gibbs' construction every effect
of the template on the ions near it is an excess per area. Its stationary value is the screen
composite
\begin{equation}\label{thsi:eq:composite}
\begin{aligned}
\DGhet={}&f\Big(\theta\big[\gslz+\dga(\DGseq,\cb,\Gs,\ds),\ \gsnz,\ \dgw\big]\Big)\\
&\times\frac{16\pi\gnl^{3}\vm^{2}}{3(\kB T\ln\Omb)^{2}} \un{J}.
\end{aligned}
\end{equation}

\paragraph{Statement.} In Eq.~\eqref{thsi:eq:composite} every effect of the template on the barrier is
counted exactly once.

\paragraph{Proof.} (1) \emph{Exactness.} Eq.~\eqref{thsi:eq:composite} is the stationary value of the
exact grand potential~\eqref{thsi:eq:grand} after the factorisation of
Section~\ref{thsi:sec:nucleation}, with $\gsl=\gslz+\dga$. (2) \emph{Channel separation.} $\Omb$ depends
on bulk activities only, so $\partial\Om/\partial(\text{any template input})=0$; all template
dependence is in $\gsl$ and $\gsn$, hence in $\theta$. $\gnl$, $\Vm$ and $T$ are template-independent
and may appear in both factors. (3) \emph{Single entry of the binding.} $\DGseq$, $\Gs$ and $\ds$ enter
$\gsl$ through Eq.~\eqref{thsi:eq:gibbs} only, and only when $\gsl$ is built from $\gslz$. They do not
enter $\gsn$, supplied from contact data with its own excess inside, and they do not enter $\Om$
(Eq.~\eqref{thsi:eq:dmu-uniform}). $\dgw$ enters $\gsn$ once. (4) Hence each template-dependent input
appears once, in one interfacial energy, and only through $\theta$. This completes the proof.

\paragraph{The two errors.} The local-ideal bookkeeping gives the enriched ion the bulk standard
state, $\Omega_\mathrm{li}=(\cl/\cb)\Omb$, so that the layer is treated as a reservoir held above the
bulk; with Ca-free interfacial energies it counts the binding once, in a channel that does not exist
at equilibrium, and with the opposite effect on the barrier. Combined with the Ca-loaded $\gsl$ it
counts the binding twice. Both exceed the exact driving force by $\kB T\ln(\cl/\cb)$ per formula unit,
$-\DGseq/\NA$ in the dilute limit, and both lower the barrier by the factor
$[\ln\Omb/\ln\Omega_\mathrm{li}]^{2}$. A second application of $\dgw$ is the third double count. The
implementation keeps these composites only as named demonstrations.

%----------------------------------------------------------------------------------------------
\subsection{Wetting branches, the spreading coefficient and the ignorance prior on $\gsn$}\label{thsi:sec:branches}

\textit{Physics.} Whether the mineral spreads on the template or beads into islands is a wetting
balance. When no measurement of the template--nucleus contact exists, its admissible range is set by
stability of a sharp interface; the range says what thermodynamics allows, not where the true value
lies.

\paragraph{Spreading.} Replacing template--liquid interface by a thick layer of mineral changes the
free energy per area by
\begin{equation}\label{thsi:eq:spread}
\Ssp=\gsl-\gsn-\gnl=\gnl(\cos\theta-1) \un{J\,m^{-2}},
\end{equation}
the second form on the branch where Young has a solution. $\Ssp\ge0$ is complete wetting (layer
growth, $f=0$); $\Ssp<0$ is island growth.

\paragraph{Bounds on $\gsn$.} Separating the template--nucleus contact into two surfaces against
vacuum costs $W_\mathrm{adh,dry}=\gp+\gamma_\mathrm{n}-\gsn$, which cannot be negative: otherwise the
solids would separate spontaneously and no contact would exist. A negative interfacial free energy
between immiscible solids would make the system create contact without limit, so $\gsn\ge0$. Hence
\begin{equation}\label{thsi:eq:gsn-bounds}
0\le\gsn\le G\equiv\gp+\gamma_\mathrm{n} \un{J\,m^{-2}},
\end{equation}
with $\gamma_\mathrm{n}$ the dry surface energy of the framing phase. The wet work
$\gsl+\gnl-\gsn$ is not bounded below: it is negative exactly on the non-wetting branch, and water can
displace a nucleus. The same two arguments give $\gow\in[0,\go+\gw]$ and $\gop\in[0,\go+\gp]$
(Section~\ref{thsi:sec:adhesion}). The ignorance prior is uniform on the admissible range.

\paragraph{Branch masses under the prior.} Young splits the $\gsn$ axis at $a=\gsl-\gnl$
(complete wetting below) and $b=\gsl+\gnl$ (non-wetting above). For $\gsn$ uniform on $[0,G]$ and the
other energies fixed,
\begin{equation}\label{thsi:eq:branches}
\begin{gathered}
P_\mathrm{complete}=\mathrm{clip}\Big(\frac{\gsl-\gnl}{G}\Big),\qquad
P_\mathrm{non}=\mathrm{clip}\Big(1-\frac{\gsl+\gnl}{G}\Big),\\
P_\mathrm{partial}=1-P_\mathrm{complete}-P_\mathrm{non} \un{1},
\end{gathered}
\end{equation}
with clip to $[0,1]$; when both cuts lie in $[0,G]$, $P_\mathrm{partial}=2\gnl/G$. Writing
$\gsl=\gp+\gw-W_\mathrm{pw}$ and $\gnl=\gamma_\mathrm{n}+\gw-W_\mathrm{nw}$, the non-wetting mass is
positive iff $W_\mathrm{pw}+W_\mathrm{nw}>2\gw$: together the two solids bind water more strongly than
water binds itself. The marginal masses are expectations over the priors of $\gsl$ and $\gnl$. They
are properties of the prior and are never reported as a probability of feasibility; flip values and
required precisions are reported instead.

\paragraph{Complete wetting is excluded on a hydrophilic template.} Because $\gsn\ge0$,
\begin{equation}\label{thsi:eq:spread-bound}
\Ssp\le\gsl-\gnl \un{J\,m^{-2}},
\end{equation}
so complete wetting needs $\gsl\ge\gnl$. Here a template is called hydrophilic when water wets it at
least as favourably, per area, as water wets the mineral: $\gsl\le\gnl$. This is not the usual
criterion of a small water contact angle in air, which by Young in air,
$\cos\theta_\mathrm{air}=(\gp-\gsl)/\gw$, constrains $\gp-\gsl$ and does not place $\gsl$ below $\gnl$.
For a template that is hydrophilic in this sense, $\Ssp\le-\gsn<0$ for every admissible $\gsn>0$:
complete wetting is excluded whatever the template--nucleus contact is, the cut point $a$ lies at or
below zero and $P_\mathrm{complete}=0$. This is a physical result, not a property of the prior.
Ca$^{2+}$ binding lowers $\gsl$ further ($\dga<0$), so the mechanism's own premise pushes the same way.
Whether a given template meets the condition is a question for the inputs, so $P(\gsl>\gnl)$, the
largest mass the spreading route can carry, is reported at every state.

\paragraph{Coalescence.} With $\Ssp<0$ a film forms only where islands meet. Equal caps of angle
$\theta$ on a square lattice of density $\Ns$ cover the cells at the cap height
$h_\mathrm{cl}=\tan(\theta/2)/\sqrt{2\Ns}$ (first contact at $\hc=h_\mathrm{cl}/\sqrt2$), so closure below
an admissible thickness $d_\mathrm{max}$ requires
\begin{equation}\label{thsi:eq:closure}
\begin{gathered}
\Ns\ge N_\mathrm{s,cl}=\frac{\tan^{2}(\theta/2)}{2d_\mathrm{max}^{2}} \un{m^{-2}},\\
d_\mathrm{max}=\max\Big(\frac{R_{A,\max}}{\rhoe},\ \frac{\ln(j_0/\jreq)}{2\kappa}\Big),\quad \kappa=\frac{\sqrt{2m_\mathrm{e}\phiB}}{\hbar} \un{m},
\end{gathered}
\end{equation}
$d_\mathrm{max}$ is the larger of the ohmic and direct-tunnelling bounds for an insulating film that
must pass the collector current; it is a limitation, not a performance figure. The lattice is the
optimistic arrangement, and the island density is taken as the upper estimate $\Nnuc$ of the next
part. A continuous film is admitted through either route: complete wetting, or closure below
$d_\mathrm{max}$.

%----------------------------------------------------------------------------------------------
\subsection{The absolute nucleation criterion (C3)}\label{thsi:sec:access}

\textit{Physics.} A barrier is low enough only relative to the attempts the system makes in the time
available. What must be achieved is set by geometry: a cathode particle touches the collector through
a contact patch, and a patch with no nucleus under it is a patch the mineral cannot cover from below.

Nuclei form on the template at the rate per area $J=J_0\e^{-\DGhet/\kB T}$, with $J_0$ the kinetic
prefactor per area. Along the registered \Rs\ states $k$, held for $\Delta t_k$, the expected areal
density is
\begin{equation}\label{thsi:eq:nnuc}
\Nnuc=\sum_kJ_0\Delta t_k\,\e^{-\Delta G^{*}_{\mathrm{het},k}/\kB T} \un{m^{-2}},
\end{equation}
an upper estimate, since it ignores exclusion zones around growing islands; a state with $\Omb\le1$
contributes nothing. If nuclei land independently, the number under one patch of area $\Apatch$ is
Poisson distributed with mean $\Nnuc\Apatch$, and the fraction of patches left empty is
$\e^{-\Nnuc\Apatch}$. Requiring that at most a fraction $q_0$ be empty,
\begin{equation}\label{thsi:eq:nreq}
\e^{-\Nnuc\Apatch}\le q_0\iff\Nnuc\ge\Nreq=\frac{\ln(1/q_0)}{\Apatch} \un{m^{-2}},\qquad q_0=1-p,
\end{equation}
with $p$ the credible level of the decision rule ($p=\NumPcred$, swept from \NumPcredSweepLow\ to
\NumPcredSweepHigh), so the admissible empty fraction moves with it. At a single state the criterion is
\begin{equation}\label{thsi:eq:bacc}
\DGhet\le B_\mathrm{acc}=\kB T\ln\frac{J_0t_\mathrm{p}}{\Nreq}=\kB T\ln\frac{J_0t_\mathrm{p}\Apatch}{\ln(1/q_0)} \un{J}.
\end{equation}
Only $\ln J_0$ enters, so an order-of-magnitude prefactor moves $B_\mathrm{acc}$ by a few $\kB T$, and
the expectation form, one nucleus per patch on average, is the special case $q_0=1/\e$.
$\Nreq$ is fixed apart from the closure density of C4, so that C3 does not borrow the island geometry
of C4.

\paragraph{The patch area.} The part of a particle touching a flat foil lies inside its projection,
and the particle touches through at least one asperity contact, so
\begin{equation}\label{thsi:eq:apatch}
a_\mathrm{asp}\le\Apatch\le\frac{\pi d_\mathrm{p}^{2}}{4},\qquad
\frac{4\ln(1/q_0)}{\pi d_\mathrm{p}^{2}}\le\Nreq\le\frac{\ln(1/q_0)}{a_\mathrm{asp}} \un{m^{2};\ m^{-2}},
\end{equation}
with $d_\mathrm{p}$ the projected-area-equivalent particle diameter. The smaller patch is the strict
end. The bound spans many decades and only its logarithm enters, so the prior is loguniform. While
either end is unsourced the criterion is reported as a gap, and the flip values of $\Apatch$ (where
the good mass equals $p$ and $1-p$) are reported in its place. The wetting part of C3 asks the
complementary question, whether the template prefers the nucleus to the liquid: $f\le\tfrac12$,
$\theta\le\pi/2$, the only value in $(0,1)$ at which the sign of the Young numerator changes.

%----------------------------------------------------------------------------------------------
\subsection{Thermodynamic stability (C5)}\label{thsi:sec:stability}

\textit{Physics.} A layer that is unstable in its environment dissolves or reacts, and the two regimes
are different environments. In the slurry the question is whether CaCO$_3$ survives alkaline water;
in the cell, whether the phases survive the collector's potential against the species the electrolyte
offers. Both are screens: necessary, not sufficient, and blind to kinetics and morphology.

\paragraph{\Rs: dissolution is decided by $\Omb$.} Carbon is held in the C(IV) carbonate system (the
regime fixes the carbonate activity, and reduction of dissolved carbonate is not observed in aqueous
processing; the construction with carbon redox allowed is a labelled sensitivity). Dissolving the layer
is then CaCO$_3$(s)$\to$Ca$^{2+}$+CO$_3^{2-}$ at the regime activities, with free energy
$\Rg T\ln\Ksp-\Rg T\ln(\aCa\aCO)=-\Rg T\ln\Omb$ per mole. The layer dissolves when $\Omb<1$ and is stable
against dissolution when $\Omb\ge1$, with the measured calcite $\Ksp$ and the speciated activities of
Section~\ref{thsi:sec:speciation}. This is the same $\Omb$ that drives nucleation: a state that cannot
nucleate CaCO$_3$ is a state in which the layer dissolves, and by Section~\ref{thsi:sec:uniform}
enrichment at the template changes neither. As a Pourbaix hull distance per non-O, non-H atom
($n_\mathrm{M}=2$), with the dissolved ions on the hull,
\begin{equation}\label{thsi:eq:pbx-omega}
\DGpbx\big|_\mathrm{diss}=\frac{\Rg T}{n_\mathrm{M}}\ln\frac{1}{\Omb}\quad(\Omb<1) \un{J\,mol^{-1}}.
\end{equation}
A computed hull therefore cannot decide dissolution. Its tolerance
$\tau=n_\mathrm{at}e_\mathrm{f}/n_\mathrm{M}$ reports dissolution only below
$\ln\Omb=-n_\mathrm{at}e_\mathrm{f}/(\Rg T)$, which hides several decades of undersaturation, and an
error $\delta$ of the computed dissolution free energy shifts the apparent $\Ksp$ by
$\delta/(\Rg T\ln10)$ decades. The test is applied at every design point of total calcium and
dissolved carbon (the dose window) and must report dissolution at an undersaturated control state
and stability at a supersaturated one; otherwise it is inconclusive.

\paragraph{\Rs: competing solids.} The Pourbaix construction
\begin{equation}\label{thsi:eq:pbx}
\begin{aligned}
\Psi_\phi={}&\frac{1}{n_\mathrm{M}}\Big[G_\phi-n_\mathrm{O}\mu_{\mathrm{H_2O}}-\Rg T\ln10\,(2n_\mathrm{O}-n_\mathrm{H})\mathrm{pH}\\
&-(2n_\mathrm{O}-n_\mathrm{H}+Q)\Fc\ESHE\Big],\\
\DGpbx={}&\Psi_\phi-\Psi_\mathrm{hull}\ge0 \un{J\,mol^{-1}},
\end{aligned}
\end{equation}
with C(IV) fixed and the dissolved Ca species removed, answers the remaining question: whether another
solid lies below CaCO$_3$, within the tolerance of computed formation energies. Adding the amorphous
offset to the edge makes the amorphous-framing test equal to the calcite-level test, because the
offset shifts $\Psi$ and leaves the hull unchanged. The native oxide's state is a baseline, not a gate.

\paragraph{\Rc: the window from explicit oxidation reactions.} Holding the collector at $\ELi$ fixes
$\muLi=\muLiz-\Fc\ELi$. A reaction that releases $\Delta n$ Li to the reservoir has
$\Delta\Phi=\Delta G_\mathrm{r}-\Delta n\Fc\ELi$ and sets the limit
\begin{equation}\label{thsi:eq:window}
\begin{gathered}
E^{*}=\frac{\Delta G_\mathrm{r}}{\Delta n\Fc} \un{V};\\
E_\mathrm{ox}=\min_{\Delta n>0}E^{*},\qquad E_\mathrm{red}=\max_{\Delta n<0}E^{*}.
\end{gathered}
\end{equation}
A Li-free phase has no Li-releasing reaction of its own, so a construction open in Li alone declares
CaCO$_3$ stable at any anodic potential by construction. Anodic stability needs reactions that hand
the carbonate's electrons to the reference and send the cation into what the electrolyte offers:
Li$^{+}$, the salt anion A$^{-}$ and the solvent. Electroneutrality pairs each released cation with A$^{-}$,
and the electrons plate Li from LiA:
\begin{equation}\label{thsi:eq:sink}
\begin{gathered}
\mathrm{CaCO_3}+2\,\mathrm{LiA}\to\mathrm{CaA_2}+\mathrm{CO_2}+\tfrac12\mathrm{O_2}+2\,\mathrm{Li},\\
\mathrm{Al_2O_3}+6\,\mathrm{LiA}\to2\,\mathrm{AlA_3}+\tfrac32\mathrm{O_2}+6\,\mathrm{Li},\\
2\,\mathrm{AlOOH}+6\,\mathrm{LiA}\to2\,\mathrm{AlA_3}+\mathrm{H_2O}+\tfrac32\mathrm{O_2}\\
+6\,\mathrm{Li},
\end{gathered}
\end{equation}
for the three gated phases (CaCO$_3$, corundum and boehmite), each with $\Delta n>0$.

\paragraph{The sink logic.} The sink is fixed by the registered electrolyte, not chosen: only its
anions are admissible, and the oxidation limit is the minimum over admissible sinks and products,
since a phase is unstable if any accessible reaction is downhill. The true product is the lower in
free energy of the solid salt and the solvated ions, so a limit computed with the solid salt is an
upper bound on the true onset. A falsification computed with it stands; a pass computed with it is
inconclusive, and a pass counts only with the solvated sink. A sink whose thermochemistry is missing
makes the limit undetermined, which never passes. The electrolyte's own anodic limit is a separate
quantity that the screen does not assess.

\paragraph{Which thermochemistry decides.} A computed reaction energy that combines database solids
with a gas reference carries the error of both, and the reference schemes bracket the measured
energies from opposite sides. Where every reactant and product is tabulated, the all-experimental
route has no such error and governs, with margin $\mexpt$: the tabulated uncertainties per transferred
electron plus a bound on any neglected solid entropy (for the registered reaction with the
perchlorate sink, $\mexpt=\NumMexpt$). A computed scheme corrected by its measured benchmark error
reproduces that same route, so it checks the bookkeeping and is never a second confirmation. Where
the experimental route is missing, the computed schemes decide with margin $m_E=n_\mathrm{unc}e_\mathrm{f}/(\Delta n\Fc)$,
counting the atoms whose computed energy does not cancel ($m_E=\NumMwindow$ for the registered
reaction): a pass must hold under the lowest-onset scheme and a falsification under every scheme. A
margin belongs to one reaction and one sink. The governing rule is
\begin{equation}\label{thsi:eq:cell-rule}
\begin{gathered}
\text{falsified if } E_\mathrm{ox}^\mathrm{expt}+\mexpt<\Ecmax\ \text{(either sink)};\\
\text{feasible only if } E_\mathrm{ox}^\mathrm{expt}-\mexpt\ge\Ecmax\\
\text{(solvated sink)},
\end{gathered}
\end{equation}
mirrored for reduction against $\Ecmin$. The charge cutoff $\Ecmax$ belongs to the cathode, so the
verdict is a function of it, with crossovers at $E_\mathrm{ox}^\mathrm{expt}\pm\mexpt$. The reading
of the limit as the onset of an observable reaction is supported by the analogous carbonate:
Li$_2$CO$_3\to2\mathrm{Li}+\mathrm{CO_2}+\tfrac12\mathrm{O_2}$ is measured to oxidise near the
experimental equilibrium limit of the same reaction, so a solid carbonate in a non-aqueous cell
oxidises close to its thermodynamic limit, with little kinetic overpotential.
